% from ega4-9-fr.tex
\oldpage[IV-3]{82}

Les parties de l'énoncé (9.7.7) concernant les propriétés (i) et (ii) se
généralisent comme suit~:

\begin{proposition}[9.7.8]
\label{IV.9.7.8-fr}
Soient $S$ un préschéma irréductible de point générique $\eta$, et $f:X\to S$
un morphisme de présentation finie. Soit $n$ (resp. $n'$) le nombre géométrique
de composantes irréductibles (resp. connexes) de $X_\eta$ (4.5.2). Alors il
existe un voisinage $U$ de $\eta$ dans $S$ tel que pour tout $s\in U$ le nombre
géométrique de composantes irréductibles (resp. connexes) de $X_s$ soit égal à
$n$ (resp. $n'$).
\end{proposition}

Compte tenu de ce que le nombre géométrique de composantes irréductibles (resp.
connexes) d'un préschéma algébrique est invariant par extension du corps de
base (4.5.6), on voit par la méthode de (9.2.3) que l'on peut se ramener au cas
où $S$ est affine, noethérien et intègre. En outre, en raisonnant comme au début
de la démonstration de (9.7.7), on voit qu'on peut supposer que les composantes
irréductibles (resp. connexes) de $X_\eta$ soient \emph{géométriquement}
irréductibles (resp. \emph{géométriquement connexes}). On sait déjà (9.7.1) que
pour tout $s$ voisin de $\eta$, $X_s$ a au moins $n$ composantes irréductibles
et $n'$ composantes connexes. D'autre part, si $Z_i$ (resp. $Z'_j$) sont les
composantes irréductibles (resp. connexes) de $X_\eta$, $X_i$ (resp. $X'_j$) le
sous-préschéma fermé réduit de $X$ ayant pour espace sous-jacent l'adhérence de
$Z_i$ (resp. $Z'_j$) dans $X$, il résulte de (9.5.1) que dans un voisinage de
$\eta$, les $(X_i)_s$ (resp. $(X'_j)_s$) forment un recouvrement de $X_s$ et que
les $(X'_j)_s$ sont deux à deux disjoints~; comme les $(X_i)_s$ (resp.
$(X'_j)_s$) sont géométriquement irréductibles (resp. géométriquement connexes)
en vertu de (9.7.7) pour $s$ voisin de $\eta$, on voit que $X_s$ a au plus $n$
composantes irréductibles et au plus $n'$ composantes connexes, d'où la
proposition.

\begin{corollary}[9.7.9]
\label{IV.9.7.9-fr}
Soit $f:X\to S$ un morphisme de présentation finie. Pour tout $s\in S$, soit
$n(s)$ (resp. $n'(s)$) le nombre géométrique de composantes irréductibles
(resp. connexes) de $f^{-1}(s)$ (4.5.2). Alors la fonction $s\mapsto n(s)$
(resp. $s\mapsto n'(s)$) est localement constructible dans $S$.
\end{corollary}

Il s'agit de montrer que la propriété «~$X$ est un préschéma algébrique sur un
corps $k$ et le nombre géométrique des composantes irréductibles (resp.
connexes) de $X$ est égal à $n$ (resp. $n'$)~» est constructible. Il résulte de
(4.5.6) que cette propriété vérifie la condition (9.2.1, (i)), et on est donc
ramené au cas où $S$ est affine, noethérien et intègre de point générique
$\eta$~; la conclusion résulte alors de (9.7.8).

\begin{corollary}[9.7.10]
\label{IV.9.7.10-fr}
Soit $X$ un préschéma localement noethérien tel que, si $X'$ est le normalisé de
$X_{\mathrm{red}}$, le morphisme canonique $f:X'\to X$ soit fini. Alors
l'ensemble des points $x\in X$ tels que $X$ soit géométriquement unibranche au
point $x$ est localement constructible dans $X$.
\end{corollary}

En effet, cet ensemble est par définition l'ensemble des points $x\in X$ tels
que le nombre de points géométriques de $f^{-1}(x)$ soit égal à 1. Mais comme
$f$ est fini, ce nombre est aussi le nombre géométrique des composantes
irréductibles de l'espace discret $f^{-1}(x)$ (compte tenu de la définition du
normalisé (\textbf{II}, 6.3.8) et de (4.5.11))~; la conclusion résulte donc de
(9.7.9).

\begin{remark}[9.7.11]
\label{IV.9.7.11-fr}
Soit $Z$ une partie localement constructible de $X$ telle que, pour tout
$s\in S$, $Z\cap X_s$ soit ouvert dans $X_s$, et désignons par $Z_s$ le
sous-préschéma de $X_s$ induit

